% Shared academic preamble. White page, black text, serif body.
\documentclass[a4paper,11pt,fontset=none]{ctexart}
\setmainfont{Noto Serif}
\setCJKmainfont{Noto Serif CJK SC}[BoldFont={Noto Serif CJK SC Bold}]
\setCJKsansfont{Noto Sans CJK SC}[BoldFont={Noto Sans CJK SC Bold}]
\setCJKmonofont{Noto Sans CJK SC}
\usepackage{geometry}
\geometry{a4paper,margin=2.3cm,headheight=14pt}
\usepackage{booktabs}
\usepackage{longtable}
\usepackage{array}
\usepackage{tabularx}
\usepackage{tikz}
\usetikzlibrary{arrows.meta,positioning}
\usepackage{float}
\usepackage{caption}
\captionsetup{font=small,labelfont=bf,skip=6pt}
\usepackage{fancyhdr}
\usepackage{setspace}
\setstretch{1.15}
\usepackage{enumitem}
\setlist[itemize]{itemsep=0.2em,topsep=0.35em,leftmargin=1.6em}
\setlist[enumerate]{itemsep=0.2em,topsep=0.35em,leftmargin=1.6em}
\usepackage{amsmath}
\usepackage{xurl}
\usepackage[hidelinks]{hyperref}
\pagestyle{fancy}
\fancyhf{}
\fancyhead[L]{\small 盼盼法式小面包}
\fancyhead[R]{}
\fancyfoot[C]{\thepage}
\renewcommand{\headrulewidth}{0.4pt}
\setcounter{secnumdepth}{3}
\setcounter{tocdepth}{2}
\emergencystretch=2em
\Urlmuskip=0mu plus 1mu
\tikzset{
  layer/.style={draw, align=flush center, font=\small, inner sep=5pt, text width=3.35cm, minimum height=1.05cm},
  wide/.style={draw, align=flush center, font=\small, inner sep=6pt, text width=14.2cm},
  flow/.style={-{Stealth[length=2.2mm]}, semithick},
  abox/.style={draw, fill=white, align=flush center, font=\scriptsize, inner sep=2.5pt,
    text width=27mm, minimum height=11mm},
  sbox/.style={draw, dashed, fill=white, align=flush center, font=\scriptsize, inner sep=2.5pt,
    text width=24mm, minimum height=11mm},
  alab/.style={font=\scriptsize\bfseries, align=center, text width=12mm},
  aflow/.style={-{Stealth[length=1.6mm]}, semithick},
  aside/.style={-{Stealth[length=1.6mm]}, densely dashed, semithick}
}
\newcommand{\contesttitlepage}[6]{%
  \begin{titlepage}
  \thispagestyle{empty}
  \centering
  \vspace*{1.5cm}
  {\large 首届江苏省“人工智能+”创新应用大赛\par}
  \vspace{0.55cm}
  {\large #1\par}
  \vspace{1.35cm}
  {\rule{0.42\textwidth}{0.5pt}\par}
  \vspace{0.75cm}
  {\Huge\bfseries #2\par}
  \vspace{0.4cm}
  {\Large #3\par}
  \vspace{0.95cm}
  {\LARGE #4\par}
  \vspace{0.5cm}
  {\large #5\par}
  \vspace{0.75cm}
  {\rule{0.42\textwidth}{0.5pt}\par}
  \vfill
  {\large #6\par}
  \vspace*{1.1cm}
  \end{titlepage}
}

\hypersetup{pdftitle={SuCompute 项目计划书},pdfauthor={盼盼法式小面包}}
\begin{document}
\contesttitlepage
  {“三苏”创新应用赛道\\赛题 AI-JIANGSU-Q36\\异构算力智能调度与能效优化}
  {SuCompute}
  {盼盼法式小面包}
  {项目计划书}
  {异构算力调度与能效的可复算策略}
  {创新组\quad 报名编号 REG-2026-000524\\命题单位：江苏省数据集团有限公司\\2026 年 11 月}

\begin{abstract}
SuCompute 在合成周上比较调度策略，并给出本地边车建议和任务凭证。种子 313、四场景等权：有效在线利用率 80.01\%，服务等级 99.83\%，相对轮询就近的单位能耗降幅 21.23\%。团队盼盼法式小面包，创新组。
\end{abstract}

\section{问题}
命题原文写明：省级调度核心加 13 个市级节点，纳管 313 个在用数据中心，总算力 271.8 EFlops，并与上海、浙江、安徽互联。调度要同时处理四件事：通算、智算、超算和国产、进口芯片的统一度量；训练、微调、推理的错配与空转；PUE、能耗和绿电进入调度；长三角多主体的服务等级与故障切换。\footnote{待核实：政府网页本次未重开。审题笔记中的平台上线日和 2026-11-18 省政府文件未进入赛题原文，不作为论据。分市座数、芯片份额和各市 PUE 不在原文中，仿真分配是假设。1/5/10 ms 为设计时延。}

\section{方案}
五分钟引擎分四层，见图~\ref{fig:arch}。数据层是 16 个节点、6 类芯片池、5 类任务，以及电价和绿电曲线。算法层用季节残差时间卷积预报未来 24 小时的到达和绿电比例（1706 个参数），用 24 期滚动整数规划做配额，再用约束策略网络在可行域内打分（1466 个参数，违约率目标 0.5\%）。服务层写出配额、故障编排和带校验和的任务凭证。交互层回放十三个设区市及上海、浙江、芜湖，并接受中文指令。

\begin{figure}[H]
\centering
\begin{tikzpicture}[x=1mm,y=1mm]
  \draw[black!40, fill=black!3] (18,71) rectangle (112,85);
  \draw[black!40, fill=black!3] (18,49) rectangle (112,63);
  \draw[black!40, fill=black!3] (18,27) rectangle (112,41);
  \draw[black!40, fill=black!3] (18,5) rectangle (112,19);
  \node[alab] at (8,78) {交互};
  \node[alab] at (8,56) {服务};
  \node[alab] at (8,34) {算法};
  \node[alab] at (8,12) {数据};
  \node[abox] at (35,78) {十三市回放\\加沪浙皖};
  \node[abox] at (67,78) {中文指令\\权重与故障};
  \node[abox] at (99,78) {边车建议\\与凭证};
  \node[abox] at (35,56) {滚动配额\\国产份额 $\ge$ 30\%};
  \node[abox] at (67,56) {可行域打分};
  \node[abox] at (99,56) {任务凭证};
  \node[abox] at (35,34) {季节残差卷积\\1706 参数};
  \node[abox] at (67,34) {24 期整数规划};
  \node[abox] at (99,34) {约束策略网络\\1466 参数};
  \node[abox] at (35,12) {16 节点\\6 类芯片};
  \node[abox] at (67,12) {合成周\\种子 313};
  \node[abox] at (99,12) {电价与绿电};
  \node[sbox] at (136,12) {小时级仿真\\单独成表};
  \node[sbox] at (136,78) {叙述模板};
  \draw[aflow] (67,19) -- (67,27);
  \draw[aflow] (67,41) -- (67,49);
  \draw[aflow] (67,63) -- (67,71);
  \draw[aside] (124,12) -- (112,12);
  \draw[aside] (124,78) -- (112,78);
  \node[font=\scriptsize, anchor=west] at (18,0) {实线：主数据流。虚线框：旁路。};
\end{tikzpicture}
\caption{SuCompute 分层架构。小时级仿真使用另一套功率系数，不与五分钟结果相减。}
\label{fig:arch}
\end{figure}

五分钟引擎峰值取 10{,}000 kW/EF，空闲系数 0.32，保活 15 分钟，碳因子用演示值 0.55 kgCO\textsubscript{2}/kWh。小时环境取 1{,}000 kW/EFlops 和分省年平均因子。全网装机在五分钟模型中为 335.8 EFlops，其中上海 26、浙江 20、芜湖 18 为假设容量。绿电抵扣碳排，不抵扣电费。本次叙述走模板，外部模型未调用。

验证集上，到达的平均绝对误差为 2.086、百分比误差 43.96\%（非零样本 45{,}179）；绿电比例的相应值为 0.0204 和 12.54\%（50{,}688）。预报训练 18 个轮次。百分比误差说明移峰仍受日周期残差限制。

\section{实测}
表~\ref{tab:kpi} 为四场景等权、种子 313 的周均值。有效在线利用率的分母含 15 分钟保活。消融、移峰次数和墙钟见佐证。

\begin{table}[ht]
\centering
\caption{同一批任务上的四条策略。}
\label{tab:kpi}
\footnotesize
\setlength{\tabcolsep}{3.5pt}
\begin{tabular}{lrrrrrr}
\toprule
策略 & 利用率 & 服务等级 & kWh/EF$\cdot$h & PUE & 绿电 & 碳强度 \\
\midrule
轮询就近 & 62.72\% & 95.72\% & 17705.0 & 1.3760 & 25.26\% & 7277.5 \\
最短队列 & 60.92\% & 98.13\% & 18181.2 & 1.3497 & 25.98\% & 7401.3 \\
成本优先 & 68.98\% & 94.97\% & 16271.5 & 1.3093 & 27.26\% & 6509.9 \\
SuCompute & 80.01\% & 99.83\% & 13945.7 & 1.2957 & 35.95\% & 4911.9 \\
\bottomrule
\end{tabular}
\end{table}

相对轮询就近：利用率 +27.57\%，服务等级 +4.12 个百分点，单位能耗 +21.23\%，碳强度 +32.51\%，PUE 降低 0.0803。脚本对设计下限印出“达到”，范围是这一合成周。本机滚动规划最长 0.0957 秒，故障编排最长 0.0629 秒。

\section{创新与应用}
同一合成周里同时优化时延圈、电价、碳、绿电和故障，并用整数规划与约束策略网络消融。每条任务有校验和凭证。边车读指标、出建议，执行仍留在现有队列。阶段 0 已可断网复跑；只读影子、人工确认和闭环是后续计划，字段对接标为待联调。

公开产品（Run:ai、Cast AI、阿里云灵骏、华为云昇腾、无问芯穹、国家超算互联网）的材料未覆盖省级多厂商配额与绿电移峰的同一张表。本队未做联合评测，也不引用市占率。SuCompute 尚未接入省平台。

没有签约和报价。\footnote{预估年电费变化 $=E\times 0.6263$ 元/kWh $\times 21.23\%$。$E$ 取 0.5 / 2.0 / 8.0 亿 kWh 时，约为 665 / 2{,}659 / 10{,}637 万元。0.6263 为小时环境引用的平段电价。精确降幅 21.232856\%。未乘碳价。}就业、客户和财政贡献未估。

团队为盼盼法式小面包。初赛材料截止 2026-11-30，12 月 1--3 日评审，12 月中旬路演。

\begin{thebibliography}{9}
\bibitem{raw} 赛题 AI-JIANGSU-Q36 原文。313 座与 271.8 EFlops 按原文引用。
\bibitem{repro} 种子 313 四场景周评测。明细表见佐证。
\end{thebibliography}
\end{document}
